\section{具身协作与学习规划方法}
\label{sec:03_10}

\subsection{观测、记忆与动作的接口}

第\ref{sec:03_09}节的联合优化以明确参数和可行域为基础。具身协作还需要回答：主体实际看到了什么，如何保留执行历史，以及规划动作能否由当前设备完成。为比较不同方法，下面建立本书的通用接口表达；它不是所引论文共有的算法或定理。

令 $\mathcal R$ 为参与协作的机器人集合，$r\in\mathcal R$ 为主体，$k\in\mathbb N_0$ 为决策事件编号。主体在该事件收到的观测为 $o_{rk}\in\mathcal O_r$，记忆为 $M_{rk}\in\mathcal M_r$，能力记录为 $C_{rk}\in\mathcal C_r$。三个集合分别表示观测、记忆和能力记录的允许取值。观测包含传感信息及已到达消息；记忆保存任务进度、已执行动作和失败原因；能力记录描述传感器、运动机构及当前健康状态。它们都是结构化对象，不具有统一单位，其中位置、时间和能量字段分别沿用 $\mathrm{m}$、$\mathrm{s}$ 和 $\mathrm{Wh}$，资格标记为无量纲布尔量。

令 $H_{rk}$ 为该主体截至事件 $k$ 实际收到的信息历史，$\mathcal A_r$ 为其技能及控制参数组成的动作空间，$\mathcal A_r^{\mathrm{info}}(H_{rk},C_{rk})\subseteq\mathcal A_r$ 为按已知信息筛选的候选动作集合。给定任务描述 $T$，策略产生动作提议
\begin{equation}
\label{s0310:policy_interface}
a_{rk}^{\mathrm{prop}}\sim
\pi_r(\cdot\mid o_{rk},M_{rk},C_{rk},T),
\qquad
\operatorname{supp}\pi_r\subseteq
\mathcal A_r^{\mathrm{info}}(H_{rk},C_{rk}),
\quad r\in\mathcal R.
\end{equation}
这里 $\pi_r$ 为动作概率分布，确定性选择是其特例；离散技能名称无量纲，连续参数采用对应控制量的单位。记忆和能力更新只能使用已收到的信息，因此相同信息历史应满足式\eqref{s0309:nonanticipativity}的非预知性要求。

候选集合表达的是信息条件下的判断。传感延迟可能使其中仍含物理上已不可执行的动作，而且各主体分别合法也不保证共享平台或联合轨迹合法。因此，动作提议还须经过执行前检查，执行后再形成反馈。令 $f_{rk}\in\mathcal F_r^{\mathrm{fb}}$ 为下一决策前收到的反馈，$\mathcal F_r^{\mathrm{fb}}$ 为反馈记录空间，则记忆更新写为
\begin{equation}
\label{s0310:memory_update}
M_{r,k+1}=\Psi_r(M_{rk},a_{rk}^{\mathrm{prop}},f_{rk}),
\qquad r\in\mathcal R,
\end{equation}
其中 $\Psi_r$ 为记忆更新规则。反馈应区分提议被拒绝、动作已执行、部分完成和完成验收，并附带相关观测及接收时刻。讨论中提出一个动作，不能直接把其预期后果写成已发生状态。

\subsection{能力描述、反馈规划与协作评价}

EMOS研究不同具身能力如何进入任务分解与分配\cite{EMOS2025}。其Robot Resume由URDF骨架的语言概括和正运动学计算共同生成，包含移动、感知和操作能力的文字及数值描述。系统接收任务、场景上下文和机器人能力档案，通过集中讨论形成子任务分配，再由各主体调用专属技能并行执行，以动作响应更新历史。能力档案因而连接任务要求与具体机构，而不是仅给机器人附加角色名称。

这一接口依赖论文明确采用的理想语义与几何重建：初始几何观测来自完美多机器人SLAM假设。其基准关闭了集成的PyBullet接触物理，并将接触到的抓取物吸附到夹爪。由此得到的任务执行结果不能证明巡检飞行、电池或接触动力学正确；机器人几何模型也不能替代本书的功率与续航模型。

COHERENT研究长时域异构任务中的提议、执行、反馈与调整，即PEFA循环\cite{COHERENT2025}。集中分配器输入任务指令、能力、各机器人局部观测及历史，输出带主体身份的子任务；执行器根据当前状态生成的动作列表选择动作，并检查其可执行性。反馈再区分任务分解错误、主体资格错误和执行受限，推动下一次提议。其局部观测不包含其他房间或关闭容器内的全部对象，因此记录历史并不等于获得完整环境状态。

该基准包含100项任务，主比较选取其中40项，覆盖5个场景，每场景8项，合计10项单类型、15项双类型和15项三类型任务。表\ref{s0310:coherent_results}仅转述论文表II的整体均值。
\begin{table}[htbp]
\centering
\caption{COHERENT论文自报的40任务比较结果}
\label{s0310:coherent_results}
\begin{tabular}{lcc}
\toprule
方法 & 成功率SR & 平均步数AS\\
\midrule
COHERENT & 0.975 & 12.4\\
移除历史 & 0.750 & 16.5\\
\bottomrule
\end{tabular}
\end{table}

令 $\mathcal Q$ 为这40项任务的集合，$G_q$ 为作者人工设计的参考完成步数，$b_q\in\{0,1\}$ 表示任务是否在 $2G_q$ 步内完成全部要求，$L_q$ 为成功时的执行步数。论文的统计规则可写为
\begin{align}
\label{s0310:coherent_metrics}
\mathrm{SR}&=\frac{1}{|\mathcal Q|}\sum_{q\in\mathcal Q}b_q,\\
\mathrm{AS}&=\frac{1}{|\mathcal Q|}\sum_{q\in\mathcal Q}\widetilde L_q,
&\widetilde L_q&=
\begin{cases}
L_q,&b_q=1,\\
2G_q+1,&b_q=0.
\end{cases}
\end{align}
成功率无量纲，步数为离散计数。失败记录保留在平均值中；完整方法的成功率仍未达到全部成功。该对照支持历史在此设置中的作用，但基准步数不是运行时间，也不包含训练、推理及通信的全部成本。这些结果不是本项目复现，不能据此作跨论文排名。

PARTNR的研究对象是家庭环境中的人机协作评价\cite{PARTNR2025}。它以自然语言任务、初始场景和主体能力约束构成问题，区分空间关系、时间顺序及只能由特定主体完成的操作。被评估规划器根据世界图和动作历史选择技能；评价函数则输入完整执行状态序列，输出完成比例、成功与否及失败解释。检查整段序列可以识别顺序错误，避免只看终态而漏掉过程要求。论文分别考察规划、感知和技能执行，报告了任务跟踪、伙伴协调及技能失败恢复方面的困难。这说明能力描述、执行反馈和独立评价承担不同功能。EMOS、COHERENT与PARTNR均不能据此被称为通用超图求解器。

\subsection{学习策略在执行链中的位置}

强化学习通过交互回报改进动作策略。HARL允许各主体使用不同参数及不同动作空间，以多主体优势分解组织贡献评估，并顺序更新策略，使后更新者考虑先更新者的变化\cite{HARL2024}。其训练输入包括交互状态、动作及共同奖励，输出异构策略；在线阶段由策略产生动作。顺序更新指训练过程，不要求机器人依次串行执行。理论推导采用完全可观测等条件，实践中的局部观测或记忆网络不能直接继承其全部保证。用于巡检还须建立训练环境、奖励和约束接口，奖励最大化也不自动等价于本书词典序目标。

模仿学习从示范中的观测与动作对应关系学习技能。NoMaD以当前及历史RGB观测、可选目标图像为输入，通过目标遮蔽统一探索与目标导航，再由扩散策略输出短时域未来动作序列，并配合拓扑记忆进行长程导航\cite{NoMaD2024}。其作用层是地面机器人的导航技能，训练依赖机器人轨迹数据。迁移时需要核对视觉条件、动作尺度和底层控制接口；已有地面导航证据不能推出多无人机调度或飞行安全保证。

OpenVLA将图像和语言指令映射为机器人操作动作，以示范数据训练动作标记预测，再将预测解码为末端位姿变化和夹爪控制\cite{OpenVLA2024}。它处于视觉语言条件下的操作策略层，适配新设备涉及动作编码、相机配置与目标域数据，并可采用微调。多种机器人上的操作能力不等于输出车队任务计划。NoMaD和OpenVLA均未在本巡检系统接入。

世界模型可在本书的接口中承担候选动作后果预测：根据当前状态估计与动作，预测后续观测或状态，供规划器比较方案。这是方法作用位置的说明，不是本项目已有实现结果。预测状态不能替代执行证据，模型误差仍须通过新观测修正。

这些方法的成本应分开计量。离线部分包括数据采集、训练、微调和验证；在线部分包括感知、记忆维护、策略推理、候选预测、通信及动作检查。扩散动作生成涉及迭代推理，语言规划涉及反复调用，预训练模型也有加载和设备适配成本。仅比较一次策略前向计算，不能代表闭环任务成本。

\subsection{故障消息到达后的合法反馈}

本书构造的教学算例C只考察热成像资格随消息接收的变化。飞机为 $U_1,U_2$，任务 $J_2$ 必须使用热成像，初始只有 $U_2$ 具备该能力。热成像故障发生于 $4\,\mathrm{min}$，规划端于 $6\,\mathrm{min}$ 收到消息；其间没有其他揭示故障的信息。令 $t\geq0$ 为以 $\mathrm{min}$ 计的时间，$\mathcal K^{\mathrm{th}}(t)$ 为实际热成像能力主体集，$\widehat{\mathcal K}^{\mathrm{th}}(t)$ 为规划端记录的能力主体集，则
\begin{align}
\label{s0310:failure_capability}
\mathcal K^{\mathrm{th}}(t)&=
\begin{cases}
\{U_2\},&0\leq t<4,\\
\varnothing,&t\geq4,
\end{cases}\\
\widehat{\mathcal K}^{\mathrm{th}}(t)&=
\begin{cases}
\{U_2\},&0\leq t<6,\\
\varnothing,&t\geq6.
\end{cases}
\end{align}
设 $a^{\mathrm{th}}(r,J_2)$ 表示将热成像任务交给主体 $r$ 的离散动作。仅按热成像资格筛选的候选动作集为
\begin{equation}
\label{s0310:failure_actions}
\widehat{\mathcal A}^{\mathrm{th}}(t)
=\{a^{\mathrm{th}}(r,J_2):r\in\widehat{\mathcal K}^{\mathrm{th}}(t)\}
=\begin{cases}
\{a^{\mathrm{th}}(U_2,J_2)\},&0\leq t<6,\\
\varnothing,&t\geq6.
\end{cases}
\end{equation}
该集合不表示平台、航迹及能量约束已经通过。在区间 $[4,6)$，信息候选仍保留 $U_2$，实际热成像动作却已不可执行；消息延迟不能被解释为动作仍然安全有效，也不能允许规划器提前使用尚未收到的故障信息。

已有Python算术重算给出的结果是：消息接收后，热成像资格集合为空，$J_2$ 仍未完成。$U_1$ 的可见光能力不能替代热成像，因此合法反馈应撤销失效资格、保留未完成状态，并在取得新的经核验能力后才允许重新分配。本例是本书定义的有限教学核对，不是公开学习算法实验。当前已有状态接续等组件验证，但完整巡检的执行与重规划闭环尚未完成；语言规划产生动作、组件执行通过与整项任务验收仍须分别判断。